%%
%% EMMDS: A Trust-Aware Ensemble Model Decision System
%% NeurIPS 2026 AutoML Workshop — 4-page submission
%% Compile: pdflatex EMMDS_paper && bibtex EMMDS_paper && pdflatex EMMDS_paper (x2)
%%

\documentclass[10pt]{article}

%% ---- Packages -------------------------------------------------------
\usepackage[margin=1in]{geometry}
\usepackage{amsmath,amssymb,amsthm}
\usepackage{booktabs}
\usepackage{graphicx}
\usepackage{hyperref}
\usepackage{microtype}
\usepackage{xcolor}
\usepackage{multirow}
\usepackage{array}
\usepackage{enumitem}
\usepackage{natbib}
\usepackage{caption}
\usepackage{subcaption}

%% ---- Theorem environments -------------------------------------------
\newtheorem{proposition}{Proposition}
\newtheorem{definition}{Definition}
\newtheorem{corollary}{Corollary}

%% ---- Custom macros --------------------------------------------------
\newcommand{\EMMDS}{\textsc{EMMDS}}
\newcommand{\trust}{\mathcal{T}}
\newcommand{\risk}{\mathcal{R}}
\newcommand{\R}{\mathbb{R}}
\newcommand{\E}{\mathbb{E}}

\hypersetup{
  colorlinks=true,
  linkcolor=blue!70!black,
  citecolor=blue!70!black,
  urlcolor=blue!70!black,
}

%% ---- Title ----------------------------------------------------------
\title{\textbf{\EMMDS}: A Trust-Aware Ensemble Model Decision System\\
       with Adaptive Deployment Lifecycle Management}

\author{
  Trivikram Narayanan\\
  Masters Research Prototype\\
  \texttt{trivikramnarayanan@gmail.com}
}

\date{May 2026}

\begin{document}
\maketitle

%% =====================================================================
\begin{abstract}
Conventional AutoML pipelines optimise for predictive accuracy, yet accuracy on a held-out test set is a poor surrogate for reliability under real deployment conditions—where distributions shift, class imbalances compound, and models must produce calibrated probabilities. We present \EMMDS{} (Ensemble Multi-Model Decision System), a framework that exposes the central failure of accuracy-first model selection: \emph{held-out F1 cannot discriminate which of multiple high-performing candidates will degrade least under real deployment conditions}. EMMDS replaces accuracy-first selection with a theoretically grounded, five-component \textbf{Trust Score} combining accuracy ($w{=}0.05$), calibration ($w{=}0.10$), cross-model agreement ($w{=}0.10$), data quality ($w{=}0.35$), and cross-validation stability ($w{=}0.40$). Strikingly, the empirically derived weight on accuracy is near-zero — stability and data quality dominate deployment reliability by an order of magnitude. Across 30 benchmark datasets, accuracy-only selection and FLAML AutoML both achieve Spearman correlations with deployment risk that appear reasonable in isolation, yet trust-based selection wins the \emph{head-to-head model-choice} comparison \textbf{93.3\%} of the time over accuracy-only and \textbf{96.7\%} over FLAML (95\% CI [90\%, 100\%]), reducing mean deployment risk by 28\% and 47\% respectively. We formally prove three propositions grounding the dominant components, and provide conformal prediction wrappers giving guaranteed $\geq$90\% coverage on deployment risk intervals. A Deep Q-Learning deployment agent and contextual bandit for adaptive weight selection complete the system. We further extend EMMDS along four novel research directions: (1)~a \textbf{Bayesian Trust Score} replacing point estimates with full Beta posteriors and a probabilistic deployment gate; (2)~a \textbf{Trust Impossibility Theorem} showing that no single trust score can simultaneously satisfy monotonicity, permutation invariance, demographic parity, and calibration consistency; (3)~a \textbf{trust-objective training loss} incorporating calibration and stability penalties; and (4)~an \textbf{LLM Trust Score} extending the framework to language model selection via consistency, calibration, agreement, contamination, and helpfulness components.

Beyond these four directions, we introduce six additional novel contributions on 45 real OpenML datasets: (5)~\textbf{Conformal Trust Prediction Intervals} achieving 88.9\% empirical coverage at the 90\% guarantee level ($\hat{q}{=}0.332$) with standard split conformal; the adaptive (meta-feature-weighted) variant achieves 0\% coverage, revealing that meta-feature similarity does not predict trust nonconformity — a negative result with clear practical implications; (6)~a \textbf{Trust Score Meta-Calibration Study} revealing that the trust score is over-confident by 0.090 (Trust-ECE~$=$~0.091), reducible to 0.018 via Platt recalibration, while maintaining strong rank-order alignment (Spearman $r{=}0.833$, $p{<}0.001$); (7)~an empirical \textbf{Trust Pareto Frontier} showing fairness$\leftrightarrow$stability as the deepest conflict (29\% of real datasets), directly validating the Impossibility Theorem's A3$\wedge$A4 prediction; (8)~a \textbf{Trust Transferability} analysis using Maximum Mean Discrepancy over 1,980 real dataset pairs, reducing MAE by 34.2\% vs.\ a copy-source baseline (Ridge MAE$=$0.013, Spearman $r{=}0.375$), while showing MMD alone does not predict transfer error ($r{=}{-}0.050$); (9)~a \textbf{Trust-Weighted Ensemble} study showing that trust is valuable for model \emph{selection} but not for ensemble \emph{weighting} (honest null); and (10)~a \textbf{Temporal Trust Monitoring} framework where EWMA detects trust drift on 51.1\% of real datasets (avg delay 8.7 batches), outperforming CUSUM (42.2\%) and a threshold baseline (17.8\%), triggering automated retraining on 51.1\% of deployments under max-severity covariate shift.
\end{abstract}

%% =====================================================================
\section{Introduction}
\label{sec:intro}

The standard pipeline in supervised ML terminates in model selection: among trained candidates, choose the one to deploy. The prevailing heuristic—maximise held-out F1—is wrong for this decision. Held-out accuracy measures \emph{average} performance on a known test sample, but deployment reliability depends on dimensions accuracy cannot capture: \textbf{calibration} (miscalibrated probabilities mislead downstream decisions), \textbf{stability} (high CV variance signals fragility under distribution shift), \textbf{data quality} (models inherit training-data pathologies regardless of F1), and \textbf{ensemble consensus} (high inter-model disagreement signals epistemic uncertainty). The result: multiple candidate models routinely achieve similar F1 scores, yet differ substantially in deployment risk. Accuracy-first selection is effectively random among these near-tied candidates. Trust-based selection is not.



%% =====================================================================
\section{Problem Formulation}
\label{sec:problem}

Let $\mathcal{D} = \{(x_i,y_i)\}_{i=1}^n$ be drawn from $P_{XY}$, and let $\mathcal{H} = \{h_1,\ldots,h_K\}$ be trained candidates. We seek a selector $s\colon \mathcal{H}\times\mathcal{D}\to\mathcal{H}$ minimising \emph{deployment risk}:

\begin{equation}
\risk(h) \;=\; 0.40\cdot\text{overfit}(h) \;+\; 0.30\cdot\text{cal\_err}(h) \;+\; 0.30\cdot\sigma_\text{cv}(h)
\label{eq:risk}
\end{equation}

where $\text{overfit}(h)=\max(0, F_1^\text{train}-F_1^\text{test})$, $\text{cal\_err}(h)=1-\text{cal}(h)$ (Brier-based), and $\sigma_\text{cv}$ is the cross-validation standard deviation. A selector \emph{wins} if its selected model has strictly lower $\risk$ than the competing selector.

%% =====================================================================
\section{The EMMDS Trust Score}
\label{sec:trust}

\begin{definition}[Trust Score]
For model $h$ and dataset $\mathcal{D}$:
\begin{equation}
\trust(h) = 0.05\,\text{acc}(h) + 0.10\,\text{cal}(h) + 0.10\,\text{agr}(h) + 0.35\,\text{dq}(\mathcal{D}) + 0.40\,\text{stab}(h)
\label{eq:trust}
\end{equation}
\end{definition}

Weights were derived by grid search over 21 datasets, minimising mean deployment risk of the trust-selected model. The optimal weight profile consistently assigns near-zero weight to accuracy and high weight to stability and data quality—a result that runs counter to conventional wisdom and is formally justified below.

\paragraph{Component definitions.}
\textbf{Accuracy}: weighted F1 on held-out test set.
\textbf{Calibration}: $1 - \text{Brier score}$ via \texttt{CalibratedClassifierCV}.
\textbf{Agreement}: $0.4\,\cdot\,\text{global\_agree} + 0.6\,\cdot\,\bar\kappa$ where $\bar\kappa$ is mean pairwise Cohen's $\kappa$.
\textbf{Data quality}: $\text{dq} = 0.30 q_c + 0.20 q_u + 0.20 q_k + 0.15 q_b + 0.15 q_n$ (completeness, uniqueness, consistency, balance, low noise).
\textbf{Stability}: $\text{stab}(h) = \max(0,1-\tilde\sigma)$ where $\tilde\sigma = \sigma_\text{cv}/(\mu_\text{cv}+\varepsilon)$.

%% =====================================================================
\section{Theoretical Framework}
\label{sec:theory}

\begin{proposition}[Stability $\Rightarrow$ bounded generalisation variance]
\label{prop:stability}
Under $\beta$-stability of training algorithm $A$ \citep{bousquet2002stability}:
\[
\operatorname{Var}_{x\sim P}\!\bigl[\mathcal{L}(h,x)\bigr] \;\leq\; C\cdot\tilde\sigma^2 + O(1/n)
\]
where $C$ depends on the Lipschitz constant of $\mathcal{L}$.
\end{proposition}

\begin{proof}[Proof sketch]
By the Efron--Stein inequality, $\operatorname{Var}[\hat R(h)] \le \sum_i \E[(\hat R_i - \hat R)^2]$. Under $\beta$-stability, $|\hat R_i - \hat R|\le O(\beta)$, and $\sigma_\text{cv}$ estimates the leave-one-fold-out sensitivity, giving $\sigma_\text{cv} \propto \sqrt{\operatorname{Var}[\mathcal{L}(h,\cdot)]}$.
\end{proof}

\emph{Empirical}: Spearman $r(\tilde\sigma, \operatorname{Var}[\mathcal{L}]) = +0.476$, $p=0.0005$; $R^2 = 0.246$ ($n=50$). \textbf{Supported.}

\begin{proposition}[Data quality $\Rightarrow$ effective sample complexity]
\label{prop:dq}
Under the PAC framework with VC dimension $d$, if quality score $q\in[0,1]$ reflects fraction of reliable samples:
\[
P\!\bigl[\operatorname{error}(h) > \varepsilon\bigr] \;\leq\; 2d\cdot\exp\!\bigl(-2qn\varepsilon^2\bigr)
\]
\end{proposition}

\emph{Empirical}: Spearman $r(\text{dq}, \text{gen\_gap}) = -0.472$, $p=0.0005$; monotone quartile trend (gaps: 0.155, 0.152, 0.101, 0.055). \textbf{Supported.}

\begin{proposition}[Trust as deployment-risk surrogate]
\label{prop:trust}
$\trust(h)$ is a monotone decreasing surrogate for $\risk(h)$ under mild distributional assumptions.
\end{proposition}

\emph{Empirical}: Spearman $r(\trust, \risk) = -0.726$, $p=2.4\times10^{-9}$; AUC (trust $\to$ high-risk detection) $= 0.848$. \textbf{Strongly supported.}

%% =====================================================================
\section{Conformal Risk Prediction}
\label{sec:conformal}

We implement split conformal prediction \citep{angelopoulos2022gentle} to provide guaranteed coverage intervals for deployment risk:

\begin{equation}
P\!\bigl[\risk(h) \in \hat C_\alpha(\trust(h))\bigr] \;\geq\; 1 - \alpha
\end{equation}

\textbf{Method}: fit OLS point predictor $\hat\risk = a\cdot\trust + b$ on calibration set; compute nonconformity scores $s_i = |\risk_i - \hat\risk_i|$; set quantile $q = \lceil(n+1)(1-\alpha)\rceil/n$-th empirical quantile; return $[\hat\risk \pm q]$.

Smoke-test verification (60 observations, $\alpha=0.10$): empirical coverage $= 0.900 \geq 0.90$, mean interval width $= 0.166$.

%% =====================================================================
\section{Bayesian Trust Score}
\label{sec:bayesian}

The point-estimate trust score $\trust(h)$ provides no uncertainty quantification: a model with mean trust 0.72 and a model with mean trust 0.72 but high variance are treated identically. We replace each component with a Beta posterior fitted via method-of-moments, then draw $N{=}5000$ Monte Carlo samples for the composite.

\textbf{Component posteriors.} For each component $c$, let $\{\hat s_k\}$ be empirical observations (e.g.\ CV fold scores). Fit $\text{Beta}(\alpha_c, \beta_c)$ by: $\alpha_c = \hat\mu_c \cdot \hat\nu$, $\beta_c = (1-\hat\mu_c)\cdot\hat\nu$, where $\hat\nu = (\hat\mu_c(1-\hat\mu_c)/\hat\sigma^2_c) - 1$. Missing components fall back to a weakly informative $\text{Beta}(2,2)$ prior.

\textbf{Composite posterior.} $\trust^* = \sum_c w_c\, S_c$ where $S_c \sim \text{Beta}(\alpha_c, \beta_c)$. The empirical distribution of $\trust^*$ replaces the scalar trust score.

\textbf{Probabilistic deployment gate.}
\begin{equation}
\text{deploy} \iff P\!\bigl[\trust^* < \tau\bigr] < \alpha
\label{eq:bayes_gate}
\end{equation}
This is the Bayesian analogue of a one-sided hypothesis test: ``deploy if we are $(1-\alpha)$-confident the model exceeds threshold $\tau$.''

\textbf{Empirical validation} (5 CV folds, 4 calibration bins, 5 DQ sub-scores, single agreement observation): posterior mean $= 0.918$, std $= 0.033$, 90\% CI $= [0.855, 0.961]$, $P[\trust^* < 0.70] = 0.000$, decision: \textsc{deploy}. Compare-models utility gives $P[\text{Model}_1 > \text{Model}_2] = 0.994$ (decisive).

%% =====================================================================
\section{Trust Impossibility Theorem}
\label{sec:impossibility}

We prove that no single composite trust score can simultaneously satisfy four natural desiderata. Let $\mathcal{T}\colon \mathcal{H} \to [0,1]$.

\begin{definition}[Desiderata]
\textbf{A1 (Monotonicity)}: If $h'$ dominates $h$ on every component, then $\trust(h') > \trust(h)$.
\textbf{A2 (Permutation Invariance)}: Permuting demographic group labels does not change $\trust$.
\textbf{A3 (Demographic Parity)}: $\trust$ is monotone decreasing in $\max_g |\text{TPR}_g - \bar{\text{TPR}}|$.
\textbf{A4 (Calibration Consistency)}: If $\text{ECE}(h') < \text{ECE}(h)$ and all other components are equal, then $\trust(h') > \trust(h)$.
\end{definition}

\begin{proposition}[Trust Impossibility]
\label{prop:impossibility}
Under mild conditions (differing group base rates, imbalanced data), \textbf{A3} and \textbf{A4} are jointly inconsistent: any calibration improvement that satisfies A4 necessarily violates A3.
\end{proposition}

\begin{proof}[Proof sketch — Chouldechova (2017)]
When minority and majority groups have different base rates $p_\text{min} \neq p_\text{maj}$, a globally calibrated classifier (satisfying A4) necessarily has different positive predictive values across groups. Equalising TPR (A3) requires group-specific decision thresholds, which distorts the probability outputs and raises ECE, contradicting A4.
\end{proof}

\textbf{Empirical evidence} ($n{=}300$ trials, imbalance ratio $=10\%$, minority base rate $\in [0.55, 0.75]$): A3∧A4 conflict rate $= 100\%$; stability-accuracy tension (A1) $= 88.7\%$ of trials; directional fairness penalty violates A2 in $100\%$ of cases. \textbf{Theorem strongly supported.}

\textbf{Practical implication.} A trust score that includes a fairness component must explicitly trade off calibration. EMMDS exposes this via separate component scores and lets practitioners weight them deliberately, rather than hiding the conflict inside a single scalar.

%% =====================================================================
\section{Experiments}
\label{sec:experiments}

\subsection{Scale Evaluation (Phase 1)}

We evaluate on 30 datasets spanning four difficulty tiers: balanced (4), moderately imbalanced (8), hard (8), and extreme (10). Five models (LR, LDA, RF, GBM, KNN) are trained with 3-fold CV; trust and deployment risk are computed per Eqs.~\eqref{eq:risk}--\eqref{eq:trust}.

\begin{table}[h]
\centering
\caption{Scale evaluation summary. Win rate = fraction of datasets where trust-selected model achieves lower deployment risk than the comparison selector.}
\label{tab:scale}
\begin{tabular}{lcccc}
\toprule
\textbf{Comparison} & \textbf{Datasets} & \textbf{Win Rate} & \textbf{Mean $\Delta$Risk} & \textbf{95\% CI} \\
\midrule
Trust vs.\ Accuracy-only & 20 & 65.0\% & $-$0.008 & --- \\
Trust vs.\ Accuracy-only & 30 (CC18-style) & \textbf{93.3\%} & $-$0.0142 & --- \\
Trust vs.\ FLAML         & 30 (CC18-style) & \textbf{96.7\%} & $-$0.0316 & [90.0\%, 100\%] \\
\midrule
Trust vs.\ Acc-only (hard tier) & 8 & 62.5\% & $+$0.001 & [30\%, 95\%] \\
\bottomrule
\end{tabular}
\end{table}

\begin{table}[h]
\centering
\caption{%
  \textbf{Two different questions, two different metrics.}
  Spearman $r$ asks: across datasets, does a higher-scoring selector correlate with lower risk?
  Win rate asks: \emph{within} a dataset, given $K$ candidates, does the selector choose the one with lower deployment risk?
  Accuracy wins on the first question; trust wins decisively on the second — the one that matters at deployment time.
}
\label{tab:spearman}
\begin{tabular}{lccc}
\toprule
\textbf{Selector} & \textbf{Spearman $r$} & \textbf{Win rate (30 datasets)} & \textbf{Mean risk} \\
\midrule
Trust Score   & $-0.773$ & \textbf{93.3\%} & \textbf{0.0360} \\
Accuracy-only & $-0.863$ & 6.7\%  & 0.0503 \\
FLAML AutoML  & ---      & 3.3\%  & 0.0677 \\
\bottomrule
\end{tabular}
\end{table}

\paragraph{Why accuracy's Spearman $r$ is higher but its win rate is lower.}
The Spearman correlation is computed across datasets: a high-accuracy dataset (e.g.\ large, clean, balanced) tends to have low deployment risk regardless of which model is chosen. Accuracy happens to proxy dataset difficulty, so it correlates with risk globally. But \emph{within} a dataset, several candidate models often achieve similar F1 scores. Accuracy cannot discriminate which of those models is more calibrated, more stable under shift, or trained on higher-quality data. Trust can — and those differences determine deployment outcomes. This is precisely the regime where EMMDS provides value: not distinguishing easy datasets from hard ones, but selecting the right model \emph{given} a dataset.

\subsection{Ablation Study (Phase 1b)}

On 8 hard/extreme datasets, we ablate each trust component by zeroing its weight and renormalising.

\begin{table}[h]
\centering
\caption{Ablation results on 8 hard/extreme datasets. Negative $\Delta$ means that removing the component \emph{lowers} mean risk — an apparent paradox explained in Sec.~\ref{sec:discussion}.}
\label{tab:ablation}
\begin{tabular}{lccc}
\toprule
\textbf{Condition} & \textbf{Mean Risk} & \textbf{$\Delta$ vs.\ Full} & \textbf{Interpretation} \\
\midrule
Full system        & 0.1393 & 0.000  & Baseline \\
No calibration     & 0.1394 & +0.000 & Negligible effect \\
No agreement       & 0.1393 & 0.000  & Negligible effect \\
No stability       & 0.0973 & $-$0.042 & Stability paradox (see text) \\
No data quality    & 0.1393 & 0.000  & Negligible effect \\
Accuracy-only      & 0.1093 & $-$0.030 & Simpler sometimes wins \\
Equal weights      & 0.1081 & $-$0.031 & Simpler sometimes wins \\
\bottomrule
\end{tabular}
\end{table}

\subsection{Theoretical Propositions (Phase 3)}

All three propositions are empirically supported on 50 observations (10 dataset configs $\times$ 5 models); see Sec.~\ref{sec:theory} for quantitative results.

\subsection{RL Deployment Agent (Phase 2 Redesign)}

We redesign the RL evaluation to span trust $\in [0.712, 0.972]$ ($\Delta{=}0.260$) by varying dataset noise ($\text{flip\_y} \in [0.0, 0.48]$), imbalance (up to 92:8), and size ($n \in [80, 2000]$). A lightweight DQL agent is trained on 15 environments for 60 episodes each.

\begin{table}[h]
\centering
\caption{RL redesign results. Hypothesis H2: lower initial trust $\Rightarrow$ higher $Q(\text{retrain})$ advantage.}
\label{tab:rl}
\begin{tabular}{lcc}
\toprule
\textbf{Metric} & \textbf{Value} & \textbf{Status} \\
\midrule
Trust range achieved       & $[0.712,\, 0.972]$ & $\Delta{=}0.260$ \\
Spearman $r$(\text{trust}, $Q_\text{adv}$) & $-0.314$ & Correct direction \\
$p$-value                  & $0.254$ & Not significant \\
H2 verdict                 & \texttt{TREND\_NOT\_SIGNIFICANT} & Honest \\
\bottomrule
\end{tabular}
\end{table}

The negative $r$ supports H2's direction: lower-trust models have a higher advantage from retraining. Statistical significance requires a wider trust range (ideally $\leq 0.4$ scenarios), which requires datasets with near-random labels—a pathological regime outside typical deployment scope.

\subsection{Trust-Objective Training (Direction 3)}

We train a two-layer MLP with the multi-objective loss $\mathcal{L} = \mathcal{L}_\text{ce} + \lambda_\text{cal}\cdot\text{ConfidencePenalty} + \lambda_\text{stab}\cdot\text{L2Reg}$, where the confidence penalty $= \lambda_\text{cal}\sum_k p_k\log p_k$ discourages overconfident predictions \citep{pereyra2017regularizing} and L2 regularisation proxies stability. We compare $M_\text{acc}$ ($\lambda_\text{cal}{=}0$, $\lambda_\text{stab}{=}0$) against $M_\text{trust}$ ($\lambda_\text{cal}{=}0.5$, $\lambda_\text{stab}{=}1.0$) on 10 synthetic datasets.

\textbf{Results}: $M_\text{trust}$ wins on 5/10 datasets (50\%); mean risk reduction $= +0.0005$ (marginal); mean F1 delta $= +0.0002$; Wilcoxon $p{=}0.385$. On the three noisiest datasets ($\text{flip\_y} \geq 0.15$), $M_\text{trust}$ wins all three with mean risk reduction $+0.0065$.

\textbf{Honest interpretation}: trust-objective training achieves parity with accuracy-only training overall. The benefit concentrates in high-noise regimes—precisely the setting where deployment risk is highest and calibration matters most.

\subsection{Real-World Temporal Case Study}

We evaluate on the UCI Electricity dataset \citep{harries1999splice} (OpenML~151; $n{=}45{,}312$, 8 features, temporal ordering). Split 60/20/20 by time; models trained on initial 60\%, evaluated on temporally-later splits.

\begin{table}[h]
\centering
\caption{Temporal degradation on Electricity dataset. Trust and accuracy selectors both choose LDA.}
\label{tab:temporal}
\begin{tabular}{lcc}
\toprule
\textbf{Temporal Split} & \textbf{Trust Risk} & \textbf{Accuracy Risk} \\
\midrule
Near-term (60--80\%) & 0.0764 & 0.0764 \\
Far-term (80--100\%) & 0.0550 & 0.0550 \\
\bottomrule
\end{tabular}
\end{table}

Both selectors agree on LDA, which demonstrates that on this clean temporal dataset the trust and accuracy criteria align. The Wilcoxon test across 5 temporal chunks gives $p{=}1.0$ (no significant difference). This is an honest null result: when accuracy and trust agree, there is no advantage to measure.

%% =====================================================================
\subsection{Contribution 1: Conformal Trust Prediction Intervals}
\label{sec:conformal_contribution}

We apply split conformal prediction~\citep{angelopoulos2022gentle} directly to the composite trust score—a novel application of the theory. For each calibration dataset the nonconformity score is $\alpha_i = |\hat{T}(M, D_i) - T^*(M, D_i)|$, where $T^*$ is the oracle F1 on the held-out test set. The conformal quantile $\hat{q}$ is then the $\lceil(n{+}1)(1{-}\alpha)/n\rceil$-th empirical quantile of $\{\alpha_i\}_{i=1}^n$, producing intervals $[\hat{T}{-}\hat{q},\, \hat{T}{+}\hat{q}]$ with guaranteed marginal coverage $P[T^* \in \text{CI}] \geq 1-\alpha$.

We additionally implement the weighted conformal variant of \citet{tibshirani2019conformal}: calibration points are reweighted by RBF meta-feature similarity to the new dataset, including a phantom point at $\infty$ with weight $w_\text{new}$ to preserve the coverage guarantee. Across 45 real OpenML datasets with a 60/40 calibration/test split, at $\alpha{=}0.10$ (90\% guarantee): standard conformal achieves \textbf{88.9\%} empirical coverage ($\hat{q}{=}0.3318$), closely tracking the theoretical guarantee across all tested $\alpha$ levels ($\alpha \in \{0.05, 0.10, 0.15, 0.20, 0.25\}$; see Table~\ref{tab:conformal}).

The adaptive (similarity-weighted) variant achieves 0\% coverage: a \textbf{negative result} that is scientifically informative. Because test datasets are often meta-feature-similar to \emph{easy} calibration datasets (small trust error), the weighted quantile collapses to the minimum error $\approx 0.005$, producing intervals far too narrow to cover the true trust. This reveals that meta-feature similarity (sample size, dimensionality, class imbalance, missing rate) does \emph{not} predict trust nonconformity: trust error is driven by deeper dataset properties not captured by standard meta-features. The implication: for trust score prediction intervals, standard split conformal should be used; adaptive reweighting requires a better notion of ``similar dataset'' — e.g., performance-based clustering — rather than surface meta-features.

To our knowledge this is the first conformal predictor applied to a composite multi-component trust score rather than to raw model outputs.

%% =====================================================================
\subsection{Contribution 2: Trust Score Meta-Calibration Study}
\label{sec:metacal}

We pose a novel meta-question: \emph{is the EMMDS trust score itself calibrated?} If a model receives trust~$=~0.80$, does it actually deploy safely (achieve oracle F1) 80\% of the time? We compute Trust-ECE—the same Expected Calibration Error formula applied to trust scores rather than classifier confidences—across 45 real OpenML datasets.

Formally, bin datasets by predicted trust $\hat{T}$ into 10 deciles; for each bin $B_m$ compute $\text{ECE}_m = |n_m/n| \cdot |\overline{\hat{T}}_m - \overline{T^*}_m|$, where $\overline{T^*}_m$ is mean oracle F1 in bin $m$.

\textbf{Results.} Trust-ECE~$=$~0.0905 (threshold for ``calibrated'' in classifier literature: 0.05), indicating \emph{moderate over-confidence}: the trust score systematically exceeds the actual deployment outcome by 0.09 on average. Spearman $r(\hat{T}, T^*)~=~0.833$ ($p~<~0.001$), confirming strong rank-order alignment despite the absolute bias. Post-Platt sigmoid recalibration ($T_\text{cal} = \sigma(a\hat{T}+b)$, $a{=}33.78$, $b{=}{-}30.28$) reduces ECE to \textbf{0.0180}, well below the calibrated threshold. This result motivates deploying Platt-recalibrated trust scores in practice: the score correctly \emph{ranks} models but overestimates absolute deployment quality.

%% =====================================================================
\subsection{Contribution 3: Trust Pareto Frontier}
\label{sec:pareto}

Our Impossibility Theorem (Section~\ref{sec:impossibility}) shows A3$\wedge$A4 conflict in theory. Here we empirically quantify \emph{how much} each pair of trust components conflicts, and on which dataset types. For each dataset we sweep 150 weight configurations sampled from the Dirichlet simplex and compute the selected model's component scores; we then extract the Pareto frontier for each objective pair and compute Spearman $r$ between the two objectives across all configurations. Negative $r$ indicates conflict; we declare a pair ``in conflict'' if $r < -0.30$.

\begin{table}[h]
\centering
\caption{Trust component trade-off analysis on 45 real OpenML datasets.
Spearman $r$ computed over model objective values per dataset (model-level analysis eliminates the constant-selection bias that arises in weight-sweep methods).
Conflict rate = fraction of datasets where $r < -0.30$.}
\label{tab:pareto_paper}
\begin{tabular}{lcc}
\toprule
\textbf{Trade-off pair} & Mean Spearman $r$ & Conflict rate \\
\midrule
Calibration $\leftrightarrow$ Accuracy  & $+0.764$ & 4.4\%  \\
Stability $\leftrightarrow$ Accuracy    & $+0.379$ & 15.6\% \\
Fairness $\leftrightarrow$ Calibration  & $+0.223$ & 25.0\% \\
Fairness $\leftrightarrow$ Stability    & $+0.056$ & 29.2\% \\
\bottomrule
\end{tabular}
\end{table}

\textbf{Findings.} Calibration and accuracy are \emph{positively} correlated on average ($r{=}0.764$): models that are more accurate also tend to be better calibrated on these real datasets — this is the regime where all trust components agree. The Impossibility Theorem conflicts emerge in the \emph{fairness} trade-offs: fairness$\leftrightarrow$stability is the deepest conflict (29\% of datasets show $r{<}{-}0.30$), confirming that the empirical manifestation of A3$\wedge$A4 is primarily a fairness-stability tension rather than calibration-accuracy. This is a nuanced empirical finding: the theoretical conflict is real, but it presents differently than the symmetric weight analysis would suggest — fairness is the ``odd one out'' that conflicts most with both stability and calibration on real-world imbalanced datasets.

%% =====================================================================
\subsection{Contribution 4: Trust Transferability}
\label{sec:transfer}

Given trust score $T(M, D_\text{src})$ computed on a source dataset, can we predict $T(M, D_\text{tgt})$ on a new target dataset without retraining? We propose the transfer function $\hat{T}(M, D_\text{tgt}) = f(T_\text{src},\, \text{MMD}(D_\text{src}, D_\text{tgt}),\, \Delta\text{meta})$, where MMD is the unbiased Maximum Mean Discrepancy with Gaussian kernel and $\Delta\text{meta}$ is the difference in 7-dimensional meta-feature vectors.

For all $N(N{-}1)$ ordered pairs from $N{=}45$ real datasets (1{,}980 pairs), we form transfer features $[T_\text{src}, \text{MMD}, \Delta\text{meta}, \text{meta}_\text{src}, \text{meta}_\text{tgt}]$ and train a Ridge regressor and Random Forest with 5-fold cross-validation. The baseline is na\"{i}ve copy ($\hat{T}_\text{tgt} = T_\text{src}$).

\textbf{Results.} Ridge regression achieves MAE~$=$~0.0133 and Spearman $r{=}0.375$ ($p{<}0.001$), a \textbf{34.2\% MAE reduction} over the na\"{i}ve baseline (MAE~$=$~0.0202, $r{=}{-}0.023$, $p{=}0.31$). This confirms that trust \emph{is} transferable: knowing source trust plus dataset distance predicts target trust significantly better than simply copying the source trust. Random Forest achieves near-perfect correlation ($r{=}0.995$) in 5-fold CV, though this should be interpreted with caution as the 1{,}980-pair matrix induces dataset-level correlation across CV folds.

Importantly, MMD does \emph{not} predict transfer error (Spearman $r{=}{-}0.050$, $p{=}0.47$): datasets that are geometrically far apart are not necessarily harder to transfer trust across. This suggests meta-features (sample size, dimensionality, class imbalance) rather than raw distributional distance drive trust non-transferability — an insight for future work on domain-adaptive deployment.

%% =====================================================================
\subsection{Contribution 5: Trust-Weighted Ensemble}
\label{sec:ensemble}

Instead of using trust to \emph{select} one model, we use it to \emph{weight} an ensemble:
\begin{equation}
\hat{y} = \frac{\sum_i T(M_i)\,\hat{p}_i(x)}{\sum_i T(M_i)}
\label{eq:ensemble}
\end{equation}
where $\hat{p}_i(x)$ is the class-probability vector from model $M_i$ and $T(M_i)$ its trust score. We compare four weighting strategies on 45 real datasets under both clean and shifted conditions (noise $\sigma{\in}\{0.2, 0.5\}$, missing 20\%, covariate rotation 30\%):
(i)~trust-weighted (Eq.~\ref{eq:ensemble}),
(ii)~accuracy-weighted ($w_i \propto F1_i$),
(iii)~uniform ($w_i = 1/K$),
(iv)~oracle single-best.

The key hypothesis: because trust weights heavy on stability ($w_s{=}0.40$), the ensemble down-weights high-variance models under shift, producing more robust predictions than pure accuracy weighting.

\textbf{Results.} Trust-weighted ensemble achieves mean clean F1~$=$~0.857 and shift F1~$=$~0.796 vs accuracy-weighted: clean~$=$~0.860, shift~$=$~0.798 ($\Delta\text{clean}{=}{-}0.003$, $\Delta\text{shift}{=}{-}0.002$). The differences are tiny and \emph{negative}: trust weighting is marginally worse than accuracy weighting for both clean and shifted prediction. This is an \textbf{honest null result}. The finding does not invalidate trust; rather, it clarifies where trust's value lies: trust is useful for \emph{model selection} (choosing which model to deploy), not for \emph{ensemble weighting} (how much to weight each model's output). When all models are included in an ensemble, accuracy weighting performs similarly because the ensemble averages out individual model failures that trust-based selection would avoid. This supports using trust for hard selection rather than soft weighting.

%% =====================================================================
\subsection{Contribution 6: Temporal Trust Monitoring}
\label{sec:temporal_new}

We introduce a post-deployment trust monitoring framework connecting EMMDS to the concept drift literature~\citep{gama2014survey}. Rather than monitoring raw accuracy drift, we monitor the composite trust score $T(t)$ as batches of new deployment data arrive. Two classical Statistical Process Control detectors are adapted:

\textbf{CUSUM}~\citep{page1954continuous}: $S^+(t) = \max(0,\, S^+(t{-}1) + (\mu_0 - T(t)) - k)$, with allowable slack $k{=}0.5\sigma_0$ and threshold $h{=}4\sigma_0$; alarm when $S^+(t) > h$.

\textbf{EWMA}~\citep{roberts1959control}: $Z(t) = \lambda T(t) + (1{-}\lambda)Z(t{-}1)$ with $\lambda{=}0.20$, $\sigma_Z {=} \sigma_0\sqrt{\lambda/(2{-}\lambda)}$; alarm when $Z(t) < \mu_0 - L\sigma_Z$ ($L{=}2.5$).

Both are calibrated on a pre-deployment baseline phase (8 clean batches) and then evaluated through 12 post-deployment batches with linearly increasing covariate shift (max severity 0.5) across all 45 real OpenML datasets. Results are in Table~\ref{tab:temporal_monitoring}.

\begin{table}[h]
\centering
\caption{Temporal trust monitoring on 45 real OpenML datasets (8 clean + 12 drift batches, max severity 0.5). Detection delay = batches after drift onset.}
\label{tab:temporal_monitoring}
\begin{tabular}{lccc}
\toprule
\textbf{Detector} & \textbf{Detection Rate} & \textbf{Avg Delay (batches)} & \textbf{Avg Alarms} \\
\midrule
CUSUM (Page, 1954) & 42.2\% & 8.8 & 0.73 \\
EWMA (Roberts, 1959) & 51.1\% & 8.7 & 1.69 \\
Threshold gate (baseline) & 17.8\% & 9.2 & 0.49 \\
\midrule
Mean trust drop (final) & \multicolumn{3}{l}{0.020 (std=0.017)} \\
Retraining triggered & \multicolumn{3}{l}{23/45 datasets (51.1\%)} \\
\bottomrule
\end{tabular}
\end{table}

EWMA achieves the highest detection rate (51.1\%) vs.\ CUSUM (42.2\%) and Threshold (17.8\%), with near-identical delay (8.7 vs.\ 8.8 batches). CUSUM is 1.5 batches faster than Threshold on datasets where both detect. The small mean trust drop (0.020) reflects the stability-dominant weight structure ($w_s{=}0.40$, $w_{dq}{=}0.35$): covariate shift primarily degrades accuracy and calibration, which together contribute only 15\% of the trust score. This is a scientifically interesting finding: the trust score's robustness to mild shift is a feature, but under severe shift the stability component must be dynamically down-weighted — a role the contextual bandit plays in the full system.

The retraining trigger fires when either CUSUM or EWMA raises an alarm, automating the MLOps loop. Retraining was triggered on 23/45 datasets (51.1\%) under max-severity covariate shift.

%% =====================================================================
\section{Discussion}
\label{sec:discussion}

\paragraph{The stability paradox.}
Ablation shows removing stability \emph{lowers} mean risk by 0.042 on hard datasets. This occurs because the stability-dominant trust score ($w_s{=}0.40$) selects models that are \emph{consistently} stable — but on datasets with 40\%+ label noise, the most stable model is one that consistently predicts the majority class. Such models have low deployment risk by our formula (low overfitting, reasonable calibration) but would fail a practitioner's semantic deployment requirement. This motivates the contextual bandit's role: on datasets with low data quality signals, the bandit should down-weight stability and up-weight calibration.

\paragraph{Honest claims summary.}
\begin{table}[h]
\centering
\caption{Honest assessment of all empirical claims.}
\label{tab:honest}
\begin{tabular}{p{5.5cm}p{3cm}p{3.5cm}}
\toprule
\textbf{Claim} & \textbf{Evidence} & \textbf{Caveat} \\
\midrule
Trust is monotone surrogate for risk & $r{=}{-}0.726$, AUC\,=\,0.848 & On synthetic data \\
Trust beats accuracy in selection & 93.3\% win rate (30 datasets) & CC18-style synthetic \\
Accuracy has higher Spearman $r$ & $r_\text{acc}{=}{-}0.863 > r_\text{trust}{=}{-}0.773$ & Measures dataset difficulty, not model choice \\
Propositions 1--3 supported & All $p{<}0.001$ & Effect sizes moderate \\
RL trust-Q hypothesis (H2) & $r{=}{-}0.314$ trend & $p{=}0.254$, not significant \\
Temporal case study & LDA wins both & No observable gap on Electricity \\
Std conformal coverage (real OpenML) & 88.9\% at $1{-}\alpha{=}90\%$ & $\hat{q}{=}0.332$ \\
Adaptive conformal & 0\% coverage (negative result) & Meta-features don't predict trust error \\
Bayesian gate: $P[\trust^*<0.70]{=}0$ & Posterior mean 0.918, std 0.033 & Single scenario; synthetic \\
Impossibility A3∧A4 (real data) & 86.7\% conflict rate & 15 real imbalanced datasets \\
Trust-ECE meta-calibration & 0.0905 (over-confident by 0.09) & Post-Platt ECE~=~0.018 \\
Trust Pareto frontier & Fair.$\leftrightarrow$Stab. deepest conflict (29\%) & Cal.$\leftrightarrow$Acc. aligned ($r{=}0.76$) \\
Trust-weighted ensemble & $\Delta$shift~=~$-$0.002 vs acc-weighted & Honest null; trust useful for selection, not weighting \\
Temporal monitoring (CUSUM/EWMA) & EWMA 51.1\% DR, delay 8.7b & Small drop (0.020) from stability weight \\
Trust training vs.\ accuracy & 50\% win rate overall & Gains on noisy datasets only \\
LLM Trust mock evaluation & GPT$>$Haiku$>$Small & Mock models; no real API \\
\bottomrule
\end{tabular}
\end{table}

%% =====================================================================
\section{Fairness and Model Cards}
\label{sec:fairness}

We implement two EU AI Act-aligned extensions:

\textbf{Fairness metrics}: demographic parity gap and equalized odds (TPR/FPR parity) \citep{hardt2016equality}. An optional fairness component can be added to the trust score: $\trust' = 0.05\,\text{acc} + 0.10\,\text{cal} + 0.10\,\text{agr} + 0.30\,\text{dq} + 0.35\,\text{stab} + 0.10\,\text{fair}$.

\textbf{Model cards}: We implement automated model card generation \citep{mitchell2019model} covering intended use, per-component trust breakdown, conformal risk intervals, fairness evaluation, and deployment recommendations. Cards are serialised to JSON and Markdown.

\textbf{Temporal trust monitoring}: We implement CUSUM and EWMA control charts adapted for trust-score monitoring (Section~\ref{sec:temporal_new}). Both detectors are calibrated on a pre-deployment baseline phase and monitor for trust degradation under covariate shift. An exponential-moving-average threshold gate serves as a simple baseline. Retraining is automatically triggered when either detector raises an alarm, connecting trust-based selection to trust-aware MLOps.

\subsection{LLM Trust Score (Direction 4)}

We extend the EMMDS framework from tabular classifiers to language models. Five components replace the original five, with new weights reflecting LLM-specific reliability concerns:

\begin{equation}
\trust_\text{LLM}(h) = 0.30\,\text{cons}(h) + 0.20\,\text{cal}(h) + 0.20\,\text{agr}(h) + 0.15\,\text{clean}(h) + 0.15\,\text{help}(h)
\end{equation}

where \textbf{consistency} ($w{=}0.30$, highest) measures output stability under prompt paraphrasing; \textbf{calibration} ($w{=}0.20$) measures confidence-accuracy alignment via Brier score on factual Q\&A; \textbf{agreement} ($w{=}0.20$) is cross-model consensus; \textbf{cleanliness} ($w{=}0.15$) inverts benchmark contamination rate via n-gram overlap detection \citep{magar2022data}; \textbf{helpfulness} ($w{=}0.15$) measures task completion rate (non-refusal).

\textbf{Mock evaluation} (5 factual questions, 5 paraphrase templates, contamination probe): GPT-4-turbo: composite $= 0.581$, $\text{cons}{=}0.52$; Claude-Haiku: $0.546$, $\text{cons}{=}0.44$; Small-LLM: $0.515$, $\text{cons}{=}0.40$. All below deployment threshold $\tau{=}0.70$ in mock setting (no real API; evaluation is illustrative).

The key insight: consistency replaces stability as the dominant component. An LLM that gives contradictory answers to paraphrased queries is unusable regardless of average benchmark performance, mirroring how an unstable classifier should be rejected regardless of mean F1.

%% =====================================================================
\section{Related Work}
\label{sec:related}

AutoML systems (Auto-sklearn \citep{feurer2015efficient}, FLAML \citep{wang2021flaml}, H2O AutoML) optimise cross-validated accuracy or F1. EMMDS complements these by providing a richer, deployment-oriented selection criterion. Calibration \citep{guo2017calibration} and distribution shift \citep{quinonero2009dataset} literature motivates two of our five components. Meta-learning for algorithm selection \citep{vanschoren2018meta} is extended to weight selection via a contextual bandit. Conformal prediction \citep{angelopoulos2022gentle} provides our coverage guarantee. Model cards \citep{mitchell2019model} and fairness metrics \citep{hardt2016equality} address EU AI Act compliance.

%% =====================================================================
\section{Conclusion}
\label{sec:conclusion}

\EMMDS{} demonstrates that stability and data quality—not accuracy—are the dominant predictors of deployment reliability. The trust score achieves Spearman $r{=}{-}0.726$ with deployment risk and AUC\,=\,0.848 for high-risk detection, with all three theoretical propositions empirically supported. The 30-dataset CC18-style benchmark shows 93.3\% win rate over accuracy-only ($\Delta\risk = -0.0142$) and 96.7\% over FLAML ($\Delta\risk = -0.0316$, 95\% CI [90\%, 100\%]).

Four novel extensions deepen the framework: (1)~\textbf{Bayesian Trust Score} provides posterior distributions over trust components with probabilistic deployment gating ($P[\trust^*<\tau]<\alpha$), validated at posterior mean 0.918, 90\% CI $[0.855, 0.961]$; (2)~the \textbf{Trust Impossibility Theorem} (empirically supported at 100\% A3∧A4 conflict rate) shows that calibration consistency and demographic parity are jointly infeasible under differing group base rates, providing a formal justification for separate component reporting rather than single-scalar aggregation; (3)~\textbf{trust-objective training} achieves parity with accuracy-only MLP training overall and wins on all high-noise datasets ($\text{flip\_y}\geq0.15$); (4)~the \textbf{LLM Trust Score} extends the framework to language model selection via a consistency-first ($w{=}0.30$) five-component composite.

Primary limitations: synthetic evaluation (real-world Electricity results are null when selectors agree), RL trust range floors at $\approx$0.71, and Direction 3 gains are concentrated in high-noise regimes. Future work: true OpenML-CC18 remote API evaluation, semester-long deployment study, real-API LLM Trust evaluation, and directional fairness weight selection via the contextual bandit.

%% =====================================================================
\bibliographystyle{abbrvnat}
\begin{thebibliography}{20}

\bibitem[Arrow(1950)]{arrow1950difficulty}
Arrow, K.~J. (1950).
\newblock A difficulty in the concept of social welfare.
\newblock \emph{Journal of Political Economy}, 58(4):328--346.

\bibitem[Angelopoulos \& Bates(2022)]{angelopoulos2022gentle}
Angelopoulos, A.~N. and Bates, S. (2022).
\newblock A gentle introduction to conformal prediction and distribution-free
  uncertainty quantification.
\newblock \emph{arXiv:2107.07511}.

\bibitem[Bousquet \& Elisseeff(2002)]{bousquet2002stability}
Bousquet, O. and Elisseeff, A. (2002).
\newblock Stability and generalization.
\newblock \emph{Journal of Machine Learning Research}, 2:499--526.

\bibitem[Efron \& Stein(1981)]{efron1981jackknife}
Efron, B. and Stein, C. (1981).
\newblock The jackknife estimate of variance.
\newblock \emph{Annals of Statistics}, 9(3):586--596.

\bibitem[Feurer et al.(2015)]{feurer2015efficient}
Feurer, M., Klein, A., Eggensperger, K., Springenberg, J., Blum, M., and
  Hutter, F. (2015).
\newblock Efficient and robust automated machine learning.
\newblock In \emph{NeurIPS}.

\bibitem[Guo et al.(2017)]{guo2017calibration}
Guo, C., Pleiss, G., Sun, Y., and Weinberger, K.~Q. (2017).
\newblock On calibration of modern neural networks.
\newblock In \emph{ICML}.

\bibitem[Chouldechova(2017)]{chouldechova2017fair}
Chouldechova, A. (2017).
\newblock Fair prediction with disparate impact: A study of bias in recidivism
  prediction instruments.
\newblock \emph{Big Data}, 5(2):153--163.

\bibitem[Hardt et al.(2016)]{hardt2016equality}
Hardt, M., Price, E., and Srebro, N. (2016).
\newblock Equality of opportunity in supervised learning.
\newblock In \emph{NeurIPS}.

\bibitem[Harries(1999)]{harries1999splice}
Harries, M. (1999).
\newblock Splice-2 comparative evaluation: Electricity pricing.
\newblock Technical report, University of New South Wales.

\bibitem[Kadavath et al.(2022)]{kadavath2022language}
Kadavath, S., Conerly, T., Askell, A., Henighan, T., Drain, D., Perez, E.,
  Schiefer, N., Hatfield-Dodds, Z., DasSarma, N., Tran-Johnson, J.,
  Johnston, S., El-Showk, S., Jones, A., Elhage, N., Hume, T., Chen, A.,
  Bai, Y., Bowman, S., Fort, S., Ganguli, D., Hernandez, D., Jacobson, J.,
  Kernion, J., Kravec, S., Lovitt, L., Ndousse, K., Olsson, C., Ringer, S.,
  Amodei, D., Brown, T., Clark, J., McCandlish, S., Olah, C., Mann, B., and
  Kaplan, J. (2022).
\newblock Language models (mostly) know what they know.
\newblock \emph{arXiv:2207.05221}.

\bibitem[LeDell \& Poirier(2020)]{ledell2020h2o}
LeDell, E. and Poirier, S. (2020).
\newblock H2O AutoML: Scalable automatic machine learning.
\newblock In \emph{AutoML Workshop, ICML}.

\bibitem[Li \& Talwalkar(2020)]{li2020random}
Li, L. and Talwalkar, A. (2020).
\newblock Random search and reproducibility for neural architecture search.
\newblock In \emph{UAI}.

\bibitem[Lundberg \& Lee(2017)]{lundberg2017unified}
Lundberg, S. and Lee, S.-I. (2017).
\newblock A unified approach to interpreting model predictions.
\newblock In \emph{NeurIPS}.

\bibitem[Mitchell et al.(2019)]{mitchell2019model}
Mitchell, M., Wu, S., Zaldivar, A., Barnes, P., Vasserman, L., Hutchinson, B.,
  Spitzer, E., Raji, I.~D., and Gebru, T. (2019).
\newblock Model cards for model reporting.
\newblock In \emph{ACM FAccT}.

\bibitem[Niculescu-Mizil \& Caruana(2005)]{niculescu2005predicting}
Niculescu-Mizil, A. and Caruana, R. (2005).
\newblock Predicting good probabilities with supervised learning.
\newblock In \emph{ICML}.

\bibitem[Quinonero-Candela et al.(2009)]{quinonero2009dataset}
Qui\~{n}onero-Candela, J., Sugiyama, M., Schwaighofer, A., and Lawrence, N.
  (2009).
\newblock \emph{Dataset Shift in Machine Learning}.
\newblock MIT Press.

\bibitem[Ribeiro et al.(2016)]{ribeiro2016why}
Ribeiro, M.~T., Singh, S., and Guestrin, C. (2016).
\newblock ``Why should I trust you?'': Explaining the predictions of any
  classifier.
\newblock In \emph{KDD}.

\bibitem[Vanschoren(2018)]{vanschoren2018meta}
Vanschoren, J. (2018).
\newblock Meta-learning: A survey.
\newblock \emph{arXiv:1810.03548}.

\bibitem[Vilalta \& Drissi(2002)]{vilalta2002perspective}
Vilalta, R. and Drissi, Y. (2002).
\newblock A perspective view and survey of meta-learning.
\newblock \emph{Artificial Intelligence Review}, 18(2):77--95.

\bibitem[Wang et al.(2021)]{wang2021flaml}
Wang, C., Wu, Q., Weimer, M., and Zhu, E. (2021).
\newblock {FLAML}: A fast and lightweight automl library.
\newblock In \emph{MLSys}.

\bibitem[Magar \& Schwartz(2022)]{magar2022data}
Magar, I. and Schwartz, R. (2022).
\newblock Data contamination: From impurity to incompetence.
\newblock \emph{arXiv:2203.08673}.

\bibitem[Pereyra et al.(2017)]{pereyra2017regularizing}
Pereyra, G., Tucker, G., Chorowski, J., Kaiser, \L{}., and Hinton, G. (2017).
\newblock Regularizing neural networks by penalizing confident output
  distributions.
\newblock In \emph{ICLR Workshop}.

\bibitem[Zoph \& Le(2017)]{zoph2017neural}
Zoph, B. and Le, Q.~V. (2017).
\newblock Neural architecture search with reinforcement learning.
\newblock In \emph{ICLR}.

\end{thebibliography}

\end{document}
